%%%%%%%% ICML 2026 EXAMPLE LATEX SUBMISSION FILE %%%%%%%%%%%%%%%%%

\documentclass{article}

% Recommended, but optional, packages for figures and better typesetting:
\usepackage{microtype}
\usepackage{graphicx}
\usepackage{subcaption}
\usepackage{booktabs} % for professional tables

% hyperref makes hyperlinks in the resulting PDF.
% If your build breaks (sometimes temporarily if a hyperlink spans a page)
% please comment out the following usepackage line and replace
% \usepackage{icml2026} with \usepackage[nohyperref]{icml2026} above.
\usepackage{hyperref}


% Attempt to make hyperref and algorithmic work together better:
\newcommand{\theHalgorithm}{\arabic{algorithm}}

% Use the following line for the initial blind version submitted for review:
\usepackage{icml2026}

% For preprint, use
% \usepackage[preprint]{icml2026}

% If accepted, instead use the following line for the camera-ready submission:
% \usepackage[accepted]{icml2026}

\usepackage{amsmath}
\usepackage{amssymb}
\usepackage{mathtools}
\usepackage{amsthm}


% if you use cleveref..
\usepackage[capitalize,noabbrev]{cleveref}

%%%%%%%%%%%%%%%%%%%%%%%%%%%%%%%%
% THEOREMS
%%%%%%%%%%%%%%%%%%%%%%%%%%%%%%%%
\theoremstyle{plain}
\newtheorem{theorem}{Theorem}[section]
\newtheorem{proposition}[theorem]{Proposition}
\newtheorem{lemma}[theorem]{Lemma}
\newtheorem{corollary}[theorem]{Corollary}
\theoremstyle{definition}
\newtheorem{definition}[theorem]{Definition}
\newtheorem{assumption}[theorem]{Assumption}
\theoremstyle{remark}
\newtheorem{remark}[theorem]{Remark}

% Todonotes is useful during development; simply uncomment the next line
%    and comment out the line below the next line to turn off comments
%\usepackage[disable,textsize=tiny]{todonotes}
\usepackage[textsize=tiny]{todonotes}

% The \icmltitle you define below is probably too long as a header.
% Therefore, a short form for the running title is supplied here:
\icmltitlerunning{Submission and Formatting Instructions for ICML 2026}

\begin{document}

\twocolumn[
  \icmltitle{Does Drawing the Causal Graph Before Answering Help? \ Evaluating Confounder-Aware Prompting on Causal Reasoning}

  % It is OKAY to include author information, even for blind submissions: the
  % style file will automatically remove it for you unless you've provided
  % the [accepted] option to the icml2026 package.

  % List of affiliations: The first argument should be a (short) identifier you
  % will use later to specify author affiliations Academic affiliations
  % should list Department, University, City, Region, Country Industry
  % affiliations should list Company, City, Region, Country

  % You can specify symbols, otherwise they are numbered in order. Ideally, you
  % should not use this facility. Affiliations will be numbered in order of
  % appearance and this is the preferred way.
  \icmlsetsymbol{equal}{*}

  \begin{icmlauthorlist}
    \icmlauthor{Firstname1 Lastname1}{equal,yyy}
    \icmlauthor{Firstname2 Lastname2}{equal,yyy,comp}
    \icmlauthor{Firstname3 Lastname3}{comp}
    \icmlauthor{Firstname4 Lastname4}{sch}
    \icmlauthor{Firstname5 Lastname5}{yyy}
    \icmlauthor{Firstname6 Lastname6}{sch,yyy,comp}
    \icmlauthor{Firstname7 Lastname7}{comp}
    %\icmlauthor{}{sch}
    \icmlauthor{Firstname8 Lastname8}{sch}
    \icmlauthor{Firstname8 Lastname8}{yyy,comp}
    %\icmlauthor{}{sch}
    %\icmlauthor{}{sch}
  \end{icmlauthorlist}

  \icmlaffiliation{yyy}{Department of XXX, University of YYY, Location, Country}
  \icmlaffiliation{comp}{Company Name, Location, Country}
  \icmlaffiliation{sch}{School of ZZZ, Institute of WWW, Location, Country}

  \icmlcorrespondingauthor{Firstname1 Lastname1}{first1.last1@xxx.edu}
  \icmlcorrespondingauthor{Firstname2 Lastname2}{first2.last2@www.uk}

  % You may provide any keywords that you find helpful for describing your
  % paper; these are used to populate the "keywords" metadata in the PDF but
  % will not be shown in the document
  \icmlkeywords{Machine Learning, ICML}

  \vskip 0.3in
]

% this must go after the closing bracket ] following \twocolumn[ ...

% This command actually creates the footnote in the first column listing the
% affiliations and the copyright notice. The command takes one argument, which
% is text to display at the start of the footnote. The \icmlEqualContribution
% command is standard text for equal contribution. Remove it (just {}) if you
% do not need this facility.

% Use ONE of the following lines. DO NOT remove the command.
% If you have no special notice, KEEP empty braces:
\printAffiliationsAndNotice{}  % no special notice (required even if empty)
% Or, if applicable, use the standard equal contribution text:
% \printAffiliationsAndNotice{\icmlEqualContribution}

\begin{abstract}
Large Language Models (LLMs) often struggle to distinguish genuine causal relationships from mere correlations. In this paper, we evaluate whether guiding LLMs to explicitly construct structured causal models---specifically identifying confounders and causal pathways---improves their causal reasoning capabilities. We propose Confounder-Aware Prompting (CAP), a three-stage prompting pipeline comprising Pathway Analysis, Confounder Identification, and Deconfounded Judgment. Evaluating on the CLadder benchmark using DeepSeek-Chat, we find that while Direct Prompting yields a poor accuracy of 58.65\%, both Standard Chain-of-Thought (CoT) (92.48\%) and CAP (91.23\%) significantly improve performance. Although CAP does not statistically outperform standard CoT in overall accuracy, it achieves high interpretability with a 78.45\% confounder identification rate. Our findings suggest that while structural prompting aids in surfacing explicit causal graphs, standard CoT effectively saturates zero-shot causal reasoning capabilities on existing standard benchmarks.
\end{abstract}

\section{Introduction}

Causal reasoning remains a fundamental challenge for artificial intelligence. Large Language Models (LLMs) exhibit remarkable capabilities in natural language understanding, yet they often fall victim to the classic fallacy of conflating correlation with causation. Recent studies suggest that generating intermediate reasoning steps, such as Chain-of-Thought (CoT), enhances performance across various logical tasks. However, it remains underexplored whether forcing an explicitly structured causal extraction---drawing the causal graph before answering---yields superior causal reasoning.

In this work, we propose \textbf{Confounder-Aware Prompting (CAP)}, a structured causal prompting method designed to elicit explicit causal graphs. We systematically evaluate CAP against Direct Prompting and Standard CoT on the CLadder dataset using DeepSeek-Chat. 

\section{Methodology}

To isolate the benefits of explicit causal structure generation, we compare three main prompting strategies:

\begin{itemize}
    \item \textbf{Direct Prompting:} Directly asking the model if the claim is valid or invalid based on the given premise.
    \item \textbf{Standard Chain-of-Thought (CoT):} Appending ``Let's think step by step'' to elicit unstructured reasoning traces before concluding.
    \item \textbf{Confounder-Aware Prompting (CAP):} A structured three-step process:
    \begin{enumerate}
        \item \textit{Pathway Analysis}: Listing all possible causal structures that explain the correlational patterns.
        \item \textit{Alternative Structure Identification}: Identifying specific structures like common causes (confounders), mediators, or colliders.
        \item \textit{Deconfounded Judgment}: Final inference combining the discovered structures to judge validity.
    \end{enumerate}
\end{itemize}

\begin{figure}[ht]
  \vskip 0.2in
  \begin{center}
    % TODO: INSERT PIPELINE FLOWCHART HERE
    % Use: \includegraphics[width=\columnwidth]{pipeline.pdf}
    \vspace{4cm} % Placeholder height
    \caption{The Confounder-Aware Prompting (CAP) Pipeline compared with Direct and Standard CoT. (Suggesting a flowchart outlining: Input $\Rightarrow$ (Direct | CoT | CAP Steps) $\Rightarrow$ Output)}
    \label{fig:pipeline}
  \end{center}
  \vskip -0.2in
\end{figure}

\section{Experimental Setup}

\paragraph{Dataset and Model:} We utilized the CLadder benchmark, evaluating samples involving complex causal relations using \texttt{deepseek-chat}.
\paragraph{Evaluation Protocol:} We measure overall accuracy (with Bootstrap 95\% CI) and False Positive Rate (FPR), where correlational links are falsely identified as causal. For CAP, we also automatically measure the Confounder Identification Rate and Reasoning Chain Completeness.
\paragraph{Statistical Significance:} We conduct pairwise McNemar's Tests with Bonferroni correction ($\alpha < 0.0083$) for rigorous comparison.

\section{Results}

\subsection{Main Performance Metrics}

Our evaluation (Table~\ref{tab:main_results}) shows that LLMs fundamentally struggle with zero-shot validation without reasoning traces. Direct Prompting achieves only $58.65\%$ accuracy (95\% CI: 55.26\%-62.16\%) and an FPR of $11.03\%$. 

In contrast, Standard CoT achieves $92.48\%$ accuracy (95\% CI: 90.60\%-94.24\%) with an FPR of $7.27\%$. Our CAP method performs competitively, reaching $91.23\%$ accuracy (95\% CI: 89.22\%-93.11\%) and an FPR of $8.77\%$. 

\begin{table}[t]
  \caption{Overall Performance Metrics on CLadder using DeepSeek-Chat ($N=798$).}
  \label{tab:main_results}
  \vskip 0.15in
  \begin{center}
  \begin{small}
  \begin{sc}
  \begin{tabular}{lccc}
  \toprule
  Method & Acc (\%) & 95\% CI & FPR (\%) \\
  \midrule
  Direct & 58.65 & 55.26-62.16 & 11.03 \\
  Std CoT & \textbf{92.48} & 90.60-94.24 & \textbf{7.27} \\
  CAP & 91.23 & 89.22-93.11 & 8.77 \\
  CAP + FS & 89.22 & 86.97-91.35 & 8.77 \\
  \bottomrule
  \end{tabular}
  \end{sc}
  \end{small}
  \end{center}
  \vskip -0.1in
\end{table}

\subsection{Statistical Analysis}

Pairwise McNemar's tests indicate that both CAP and Standard CoT significantly outperform Direct Prompting ($p \approx 6.26 \times 10^{-46}$ and $p \approx 3.63 \times 10^{-51}$, respectively). However, the difference between CAP and Standard CoT is \textit{not} statistically significant ($p=0.2113$). Adding few-shot examples (CAP+FS) surprisingly deteriorated performance compared to zero-shot CAP, suggesting complex prompting might restrict the model's generalized logic.

\subsection{CAP Quality and Interpretability}

Although CAP does not surpass CoT in pure accuracy, it enforces structural rigor. Automatic metrics demonstrate that CAP executed the complete 3-step reasoning chain 100\% of the time, and successfully surfaced explicit confounder identifications in $78.45\%$ of instances.

\section{Conclusion}

Structure-aware prompts effectively mitigate correlation-causation fallacies, vastly outperforming Direct Prompting. However, explicitly forcing a rigid causal graph formulation (CAP) performs comparably to generic Chain-of-Thought prompting, suggesting that modern LLMs possess internal heuristics that are successfully saturated by unstructured reasoning pathways alone.

\section*{Accessibility}

Authors are kindly asked to make their submissions as accessible as possible
for everyone including people with disabilities and sensory or neurological
differences. Tips of how to achieve this and what to pay attention to will be
provided on the conference website \url{http://icml.cc/}.

\section*{Software and Data}

If a paper is accepted, we strongly encourage the publication of software and
data with the camera-ready version of the paper whenever appropriate. This can
be done by including a URL in the camera-ready copy. However, \textbf{do not}
include URLs that reveal your institution or identity in your submission for
review. Instead, provide an anonymous URL or upload the material as
``Supplementary Material'' into the OpenReview reviewing system. Note that
reviewers are not required to look at this material when writing their review.

% Acknowledgements should only appear in the accepted version.
\section*{Acknowledgements}

\textbf{Do not} include acknowledgements in the initial version of the paper
submitted for blind review.

If a paper is accepted, the final camera-ready version can (and usually should)
include acknowledgements.  Such acknowledgements should be placed at the end of
the section, in an unnumbered section that does not count towards the paper
page limit. Typically, this will include thanks to reviewers who gave useful
comments, to colleagues who contributed to the ideas, and to funding agencies
and corporate sponsors that provided financial support.

\section*{Impact Statement}

Authors are \textbf{required} to include a statement of the potential broader
impact of their work, including its ethical aspects and future societal
consequences. This statement should be in an unnumbered section at the end of
the paper (co-located with Acknowledgements -- the two may appear in either
order, but both must be before References), and does not count toward the paper
page limit. In many cases, where the ethical impacts and expected societal
implications are those that are well established when advancing the field of
Machine Learning, substantial discussion is not required, and a simple
statement such as the following will suffice:

``This paper presents work whose goal is to advance the field of Machine
Learning. There are many potential societal consequences of our work, none
which we feel must be specifically highlighted here.''

The above statement can be used verbatim in such cases, but we encourage
authors to think about whether there is content which does warrant further
discussion, as this statement will be apparent if the paper is later flagged
for ethics review.

% In the unusual situation where you want a paper to appear in the
% references without citing it in the main text, use \nocite
\nocite{langley00}

\bibliography{example_paper}
\bibliographystyle{icml2026}

%%%%%%%%%%%%%%%%%%%%%%%%%%%%%%%%%%%%%%%%%%%%%%%%%%%%%%%%%%%%%%%%%%%%%%%%%%%%%%%
%%%%%%%%%%%%%%%%%%%%%%%%%%%%%%%%%%%%%%%%%%%%%%%%%%%%%%%%%%%%%%%%%%%%%%%%%%%%%%%
% APPENDIX
%%%%%%%%%%%%%%%%%%%%%%%%%%%%%%%%%%%%%%%%%%%%%%%%%%%%%%%%%%%%%%%%%%%%%%%%%%%%%%%
%%%%%%%%%%%%%%%%%%%%%%%%%%%%%%%%%%%%%%%%%%%%%%%%%%%%%%%%%%%%%%%%%%%%%%%%%%%%%%%
\newpage
\appendix
\onecolumn
\section{You \emph{can} have an appendix here.}

You can have as much text here as you want. The main body must be at most $8$
pages long. For the final version, one more page can be added. If you want, you
can use an appendix like this one.

The $\mathtt{\backslash onecolumn}$ command above can be kept in place if you
prefer a one-column appendix, or can be removed if you prefer a two-column
appendix.  Apart from this possible change, the style (font size, spacing,
margins, page numbering, etc.) should be kept the same as the main body.
%%%%%%%%%%%%%%%%%%%%%%%%%%%%%%%%%%%%%%%%%%%%%%%%%%%%%%%%%%%%%%%%%%%%%%%%%%%%%%%
%%%%%%%%%%%%%%%%%%%%%%%%%%%%%%%%%%%%%%%%%%%%%%%%%%%%%%%%%%%%%%%%%%%%%%%%%%%%%%%

\end{document}

% This document was modified from the file originally made available by
% Pat Langley and Andrea Danyluk for ICML-2K. This version was created
% by Iain Murray in 2018, and modified by Alexandre Bouchard in
% 2019 and 2021 and by Csaba Szepesvari, Gang Niu and Sivan Sabato in 2022.
% Modified again in 2023 and 2024 by Sivan Sabato and Jonathan Scarlett.
% Previous contributors include Dan Roy, Lise Getoor and Tobias
% Scheffer, which was slightly modified from the 2010 version by
% Thorsten Joachims & Johannes Fuernkranz, slightly modified from the
% 2009 version by Kiri Wagstaff and Sam Roweis's 2008 version, which is
% slightly modified from Prasad Tadepalli's 2007 version which is a
% lightly changed version of the previous year's version by Andrew
% Moore, which was in turn edited from those of Kristian Kersting and
% Codrina Lauth. Alex Smola contributed to the algorithmic style files.
